\documentclass[10pt,a4paper]{amsart}
\usepackage[margin=19mm]{geometry}
\usepackage{amsmath,amssymb,mathtools}
\usepackage[scr=rsfs,cal=euler]{mathalfa}
\usepackage{xfrac}
\usepackage[unicode,hidelinks,bookmarks=false]{hyperref}
\newcommand{\ie}{i.e.\ }
\newcommand{\cf}{cf.\ }
\newcommand{\resp}{respectively\ }
\newcommand{\Subsection}{\subsection}
\providecommand{\nobreakdash}{\nobreak\hbox{-}}
\setlength{\emergencystretch}{2em}
\multlinegap0pt
\usepackage{bm}
\usepackage{bbm}
\usepackage{dsfont}
\usepackage{xspace}
\usepackage{cancel}
\usepackage{mathtools}
\usepackage{amsthm}
\makeatletter
\@ifundefined{thm}{\newtheorem{thm}{Theorem}}{}
\makeatother
\newcommand{\tvs}{\mathrm{tvs}}
\title[A Unified Approach to Total Vertex Irregularity Strength]{A Unified Approach to Total Vertex Irregularity Strength}
\author{Aleams Barra}
\date{Source: arXiv:2506.17922v2 (CC BY 4.0)}
\begin{document}
\maketitle

\noindent\emph{Source and licence.} A Unified Approach to Total Vertex Irregularity Strength. Version arXiv:2506.17922v2 (CC BY 4.0), distributed under the Creative Commons Attribution 4.0 International licence, as recorded in the arXiv licence metadata for this exact version and shown on the abstract page https://arxiv.org/abs/2506.17922v2. Full rights notice: licenses/arxiv-2506.17922-CC-BY-4.0.txt.\par\smallskip

\noindent\textit{Editorial scope.} Counted unit: the Thm whose source label is \texttt{None} in the article (source0.tex lines 435--440, source label \texttt{None}), delivered together with the article's own complete original proof of it (source0.tex lines 442--546, one \texttt{proof} environment of 105 lines). No mathematics is written, repaired or re-derived: statement and proof are the article's own text line for line. Required context retained uncounted: cycle (lines 113--118). Every definition, display and label of those lines is retained unchanged. Omitted from this package and not redistributed: 1--112, 119--434, 441--441, 547--557. Nothing inside a retained excerpt is omitted or altered: every retained body is delivered exactly as the article's lines are written, so no omission receipt of any kind accompanies this package. Citations: the retained excerpts carry no citation command at all, so no bibliography entry of the article is transcribed and no reference is redistributed. Reference adaptation: none was needed and none was made; every reference inside the delivered unit resolves inside it, so no unresolved reference or citation is produced. Printed slips kept literal: no printed word, symbol, hypothesis or number of the counted statement or of its proof was altered. Layout and class: the article's own class is \texttt{article} with packages including fontenc, inputenc, amsmath, amsthm, amscd, amsfonts, amssymb, graphicx, color, geometry, hyperref; the wrapper is the lane's pinned \texttt{amsart} editorial wrapper at 10pt on A4, which restates the theorem-like environments the retained text needs and adds the article's own macro definitions that the retained text needs (\texttt{\textbackslash tvs}), with extra wrapper packages (bm, bbm, dsfont, xspace, cancel, mathtools). The delivered numbering is the unit's own. Assets: no retained span contains a figure environment, a graphics-inclusion command, a PDF-page inclusion, a table or a TikZ picture; no image file is copied into this package, the package declares no asset dependency and no image file is redistributed with it. Licence: version arXiv:2506.17922v2 of this article is distributed under the Creative Commons Attribution 4.0 International licence, as recorded in the arXiv licence metadata for this exact version and shown on the abstract page https://arxiv.org/abs/2506.17922v2. A full rights notice is delivered at licenses/arxiv-2506.17922-CC-BY-4.0.txt. Characters: every retained excerpt is pure ASCII, so the delivered engine, XeLaTeX with its default Latin Modern text font, reports no missing character.
\begin{thm}\label{cycle}
Let \(C_n\) be the cycle graph with \(n\) vertices. Then
\[
\mathrm{tvs}(C_n) = \lceil (n+2)/3 \rceil.
\]
\end{thm}

\begin{thm}
Let $G$ be a simple $2$-regular graph of order $n$. Then
\[
\tvs(G)=\left \lceil \frac{n+2}{3}\right\rceil.
\]
\end{thm}

\begin{proof}
The graph $G$ is a disjoint union of cycles, $G=\bigcup_{i=1}^m C_i$, where each cycle has length at least $3$. Since the case $m=1$ corresponds to a single cycle (Theorem~\ref{cycle}), we assume $m\ge 2$, which implies $n\ge 6$. Without loss of generality, assume the cycle lengths are ordered such that $n_1 \le n_2 \le \dots \le n_m$, where $n_i = |V(C_i)|$.

Let $s=\lceil (n+2)/3\rceil$. We first consider the case where $s$ is odd. We construct a $\{1,s\}$-edge labeling such that the partition sets $V_2$, $V_{s+1}$, and $V_{2s}$ (sets of vertices with vertex sums $2$, $s+1$, and $2s$, respectively) have sizes
\[
|V_2|=s, \quad |V_{s+1}|=s-1, \quad |V_{2s}|=n+1-2s.
\]
Let $k\ge 0$ be the maximal integer such that $\sum_{i=1}^{k} n_i \le s$, where the empty sum is interpreted as $0$. We write
\[
\sum_{i=1}^{k} n_i = s-a \qquad \text{for some } a \ge 0.
\]
We construct the edge labeling in stages to fix $|V_2|$ exactly at $s$.

\smallskip
\noindent
\textbf{Case 1: $n_{k+1} \ge a+2$.} \\
Label all edges of the first $k$ cycles $C_1, \dots, C_k$ with $1$. On $C_{k+1}$, label a path of $a+1$ consecutive edges with $1$. The remaining edges of $C_{k+1}$ and all edges of subsequent cycles are initially left unlabeled. The cycles $C_1, \dots, C_k$ contribute exactly $\sum_{i=1}^k n_i = s-a$ vertices to $V_2$. The path of $a+1$ edges labeled $1$ in $C_{k+1}$ contributes exactly $(a+1)-1 = a$ internal vertices to $V_2$. Thus, the total count is $|V_2| = (s-a) + a = s$.

\smallskip
\noindent
\textbf{Case 2: $n_{k+1} = a+1$.} \\
If we were to label all $a+1$ edges of $C_{k+1}$ with $1$, the cycle would contribute $a+1$ vertices to $V_2$, yielding a total of $|V_2| = (s-a) + (a+1) = s+1$. Since replacing a label $1$ with $s$ alters the vertex sum of both endpoints, it reduces $|V_2|$ by exactly $2$, making it impossible to achieve $|V_2|=s$ using $C_{k+1}$ alone.

Instead, we label exactly one edge of $C_{k+1}$ with $s$ and the remaining $n_{k+1}-1$ edges with $1$. Then $C_{k+1}$ contributes
\[
n_{k+1}-2=a-1
\]
vertices of vertex sum $2$, so the current total is
\[
(s-a)+(a-1)=s-1.
\]
The two endpoints of the edge labeled $s$ in $C_{k+1}$ now have vertex sum $s+1$. To recover the missing vertex for $V_2$, we use the next cycle, $C_{k+2}$. Note that
\[
\sum_{i=1}^{k+1} n_i = s+1 =\left\lceil\frac{n+2}{3}\right\rceil+1 < \frac{n+2}{3}  + 2< n \quad \text{for } n \ge 6,
\]
ensuring $C_{k+2}$ exists. On $C_{k+2}$, we label exactly two consecutive edges with $1$. This creates exactly one vertex of vertex sum $2$, namely the vertex between the two edges labeled $1$. Therefore
\[
|V_2|=(s-1)+1=s.
\]
The remaining edges of $C_{k+2}$ are left unlabeled for now, which completes Case 2.

\smallskip
\noindent
Thus, in either case, we have fixed $|V_2|=s$. We now determine the labels for the remaining unlabeled edges so as to obtain $|V_{s+1}| = s-1$.

We now complete the labeling of the current cycle (if it is only partially labeled) and label every subsequent fresh cycle using alternating patterns that produce an even number of vertices with vertex sum $s+1$, while creating no new pair of consecutive edges labeled $1$. We use patterns such as $s,1,s,1,\ldots,s,1$ or $s,1,s,1,\ldots,s,1,s$ on a fresh cycle, and patterns such as $s,\mathbf{1},s,1,\ldots,s,1$ or $s,\mathbf{1},s,1,\ldots,s,1,s$ on a cycle that already contains a block of $1$s (denoted by $\mathbf{1}$). Among the \(n-s\) vertices outside \(V_2\), at most one vertex in each remaining cycle has vertex sum \(2s\). Hence the number of vertices with vertex sum \(s+1\) that is produced is at least
\[
n-s-(m-k).
\]
In any case, we will show that this quantity is at least \(s-1\). We now distinguish two subcases. If \(k=0\), then necessarily \(n_1\ge s+1\). It is impossible that \(m\ge 3\), for otherwise
\[
n=\sum_{i=1}^m n_i \ge n_1+n_2+n_3 \ge 3(s+1)>n,
\]
a contradiction. Hence \(m\le 2\). Since we are assuming \(m\ge 2\), it follows that \(m=2\). Therefore \(n=n_1+n_2\ge 2(s+1)\), and after \(|V_2|=s\) has been fixed, the number of vertices with vertex sum \(s+1\) that can be produced is at least
\[
n-s-2\ge s>s-1.
\]
If \(k\ge 1\), then \(m-k\le m-1\), and since every cycle has length at least \(3\), we have \(m\le \left\lfloor n/3 \right\rfloor\). Hence
\[
n-s-(m-k)\ge n-s-\left\lfloor \frac n3 \right\rfloor +1 \ge s-1.
\]
If the resulting number is larger than \(s-1\), we reduce it by replacing local patterns \(s,1,s\) with \(s,s,s\). This operation decreases \(|V_{s+1}|\) by exactly \(2\) and does not affect \(|V_2|\). Indeed, the only vertices of vertex sum \(s+1\) not arising from a local pattern \(s,1,s\) are the two boundary vertices adjacent to the initial block of edges labeled \(1\). Therefore, whenever \(|V_{s+1}|\ge 4\), at least two vertices of \(V_{s+1}\) arise from some local pattern \(s,1,s\), so such a replacement is possible. Since each replacement decreases \(|V_{s+1}|\) by \(2\), the last step before reaching the target \(|V_{s+1}|=s-1\) occurs when \(|V_{s+1}|=s+1\). As \(n\ge 6\), we have \(s=\left\lceil\frac{n+2}{3}\right\rceil\ge 3\), and hence \(s+1\ge 4\), so this final reduction can still be performed. Repeating this operation as needed, we obtain exactly \(|V_{s+1}|=s-1\), while \(|V_2|\) remains unchanged.

Consequently, every vertex not in \(V_2\) or \(V_{s+1}\) has vertex sum \(2s\). Thus we obtain the distribution
\[
|V_2|=s,\qquad |V_{s+1}|=s-1,\qquad |V_{2s}|=n+1-2s.
\]
This is identical to the distribution in Case~2 of Theorem~\ref{cycle}. Therefore, we assign vertex labels to the ordered sets \(\mathbf{V_2}\), \(\mathbf{V_{s+1}}\), and \(\mathbf{V_{2s}}\) according to the same pattern as in that case. Since this assignment depends only on the distribution of these sets and not on the structure of the underlying graph, it guarantees that all vertex weights are pairwise distinct and bounded by $s$.

It remains to consider the case when $s$ is even. The construction used above to force the size of $V_2$ does not rely on the parity of $s$. Hence we may run the same first-stage procedure with target
\[
|V_2|=s-1.
\]
After this stage is fixed, there remain
\[
n-(s-1)
\]
vertices outside $V_2$, distributed among the remaining $(m-k)$ cycle-components under consideration. Using the same alternating completion patterns as in the odd case, at most one vertex per such component fails to contribute to $V_{s+1}$. Therefore
\[
|V_{s+1}|\ge n-(s-1)-(m-k).
\]
Using the inequality proved above,
\[
n-s-(m-k)\ge s-1,
\]
we get
\[
|V_{s+1}|\ge n-(s-1)-(m-k)=(n-s-(m-k))+1\ge (s-1)+1=s.
\]
Hence we can first obtain $|V_{s+1}|\ge s$. The contribution to $V_{s+1}$ from each cycle is even, so the total number $|V_{s+1}|$ obtained after the alternating completion is even. If this number is larger than $s$, we repeatedly apply the local replacement
\[
(s,1,s)\longrightarrow (s,s,s),
\]
which decreases $|V_{s+1}|$ by exactly $2$ and leaves $|V_2|$ unchanged. As in the odd case, whenever a reduction is needed there are enough local patterns $s,1,s$ available outside the initial block of $1$'s to perform it. Since $s$ is even, after finitely many such steps we obtain exactly $|V_{s+1}|=s$.

Consequently, we obtain an edge labeling such that
\[
|V_2|=s-1,\quad |V_{s+1}|=s,\quad \text{and}\quad |V_{2s}|=n+1-2s.
\]
This distribution coincides with that of Case~1 in Theorem~\ref{cycle}. Although the underlying graphs are different, the final vertex-labeling step depends only on the partition \(V(G)=V_2\sqcup V_{s+1}\sqcup V_{2s}\) and the cardinalities of these three classes. Since these coincide with those in Case~1 of Theorem~\ref{cycle}, we assign vertex labels to the ordered sets \(\mathbf{V_2}\), \(\mathbf{V_{s+1}}\), and \(\mathbf{V_{2s}}\) by the same scheme, and the distinct-weight verification applies verbatim. Therefore,
\[
\tvs(G)=\left\lceil\frac{n+2}{3}\right\rceil.
\]

\end{proof}

\end{document}
